\documentclass[10pt]{article} 

\usepackage[utf8]{inputenc}

\usepackage{multicol}
\setlength{\columnsep}{1cm}

\usepackage{abc}
\usepackage{graphicx}

\usepackage[margin=0.5in]{geometry} 
\geometry{letterpaper}

\usepackage{fancyhdr} 
\pagestyle{fancy} % options: empty , plain , fancy
\renewcommand{\headrulewidth}{0pt} % customise the layout...
\lhead{}\chead{}\rhead{}
\lfoot{}\cfoot{\thepage}\rfoot{}

\begin{document}

\thispagestyle{empty}

\begin{abc}[name=jefferson,options={-O= +c -F single.fmt}]
X:1
T:Jefferson and Liberty
T:The Gobby-O
M:6/8
L:1/8
K:G
B |: "Am"c2A AGA | E2A ABc | "G"B2G GBc | dBG GAB |
"Am"c2A AGA | E2A A2e | "C"edc "G"BAG |1 "E"E2A"Am"A2B :|2 "E"E2A "Am"AcB ||
"Am"A2Bc2d | e2fg3 | "C"e2fg2e | "G"dBG G2E | 
"Am"A2Bc2d | efga3 | "C"edc "G"BAG |1 "E"E2 A "Am"AcB :|2 "E"E2A "Am"A3 ||
\end{abc}

% \vskip 0.5in
\noindent\makebox[\textwidth][c]{%
\begin{minipage}{6.5in}
  \setlength{\parindent}{1.5em}
  
  The tune known as ``Jefferson and Liberty'' is thought to have evolved from the jig, ``An Gobbi O'' (in Irish) or ``The Gobby O'' translating roughly into ``The Mouthy Boy'', which is \#105 in \textit{O’Neill’s 1001: The Dance Music of Ireland}. The lyrics were composed by ornithologist, painter, and poet, Alexander Wilson, in celebration of Thomas Jefferson’s 1801 inauguration as president, and in protest of the Alien and Sedition Acts signed into law by President John Adams in 1798. In Wilson’s notebook, the tune given is ``Vive la souverain People'' --- not a tune we have been able to find any mention of elsewhere --- and when it was published in the Aurora Magazine in January, 1801, the suggested tune was ``Willie Was a Wanton Wag``. It is not clear when the present melody picked up the title.

  Other suggested etymologies for the ``Gobby O'' title include a corruption of ``Gaberlunzie Boy'' (gaberlunzie being a medieval Scots word for a licensed beggar) and a claim that ``gobby-stick'' is a term for a fife or flute.

  Several versions of the lyrics have been published. Here is one version:
  
  \begin{multicols}{2}
    {\parindent0pt 
      \raggedright
      \small
      The gloomy night before us flies, \\
      The reign of terror now is o'er; \\
      Its gags, inquisitors, and spies, \\
      Its herds of harpies are no more! \\
\bigskip
      Chorus: \\
      Rejoice, Columbia's sons, rejoice! \\
      To tyrants never bend your knee, \\
      But join with heart, and soul, and voice, \\
      For Jefferson and Liberty! \\
\bigskip

      No lordling here, with gorging jaws \\
      Shall wring from industry the food; \\
      Nor fiery bigot's holy laws \\
      Lay waste our fields and streets in blood! \\
\bigskip
      Here strangers from a thousand shores \\
      Compelled by tyranny to roam, \\
      Shall find amidst abundant stores, \\
      A nobler and happier home. \\
\bigskip

      Here Art shall lift her laurelled head, \\
      Wealth, Industry, and Peace, divine; \\
      And where dark, pathless forests spread, \\
      Rich fields and lofty cities shine. \\
\bigskip
      From Europe's wants and woes remote, \\
      A friendly waste of waves between, \\
      Here plenty cheers the humblest cot, \\
      And smiles on every village green. \\
\bigskip
      Here free as air, expanded space, \\
      To every soul and sect shall be -- \\
      That sacred privilege of our race -- \\
      The worship of the Deity. \\
\bigskip
      Let foes to freedom dread the name; \\
      But should they touch the sacred tree, \\
      Twice fifty thousands swords would flame \\
      For Jefferson and liberty. \\
\bigskip
      From Georgia to Lake Champlain, \\
      From seas to Mississippi's shore, \\
      Ye sons of freedom loud proclaim \\
      ``The reign of terror is no more.'' \\
\bigskip
      These gifts, great Liberty, are thine, \\
      Then thousand more we owe to thee, \\
      Immortal may their mem'ries shine, \\
      Who fought and died for Liberty.
      
    }
    
  \end{multicols}
  
\end{minipage}
}


\end{document}
